\documentclass[12pt]{article}

\usepackage[a4paper, total={6in, 8in}]{geometry}
\usepackage{textcomp}

\begin{document}
\setlength{\parindent}{0pt}

\def \t {\textrightarrow}
\def \v {\vspace{1em}}
\def \bi {\begin{itemize}}
\def \ei {\end{itemize}}
\def \s[#1] {\section*{#1}}
\def \ss[#1] {\subsection*{#1}}
\def \sss[#1] {\subsubsection*{#1}}

\s[Guerra fredda]
È chiaro che non si vuole lo scontro \t c'è una provocazione in atto, che poi caratterizzera tutta la seconda

\ss[Guerra di Korea]
Altra zona di crisi oltre Berlino \t e l'asia sempra cadere in mano sovietica (anche in cina governo comunista con Mao) \\
L'asia sembrava persa per l'occidente \t ma la guerra di Korea mostra che non è cosi \\
Nel 1950 scoppia questa guerra perchè alla fine della guerra la Korea era divisa in Nord (comunista) e Sud (alleata con america) \t e sono due governi distinti \\
Entrambe le Koree rivendicano l'intero territorio nazionale pero \t e la frontiera diventa teatro di tensione \\
La korea del nord poi invade quella del sud e viene sostenuta dall URSS \t interviene l ONU che condanna korea del nord, e quindi manda truppe \\
America quindi inteviene sotto la bandiera dell'ONU \t respingono i nord koreani, oltrepassano il confine e invadono la korea del nord \t e interviene cosi la Cina \\
La cina ha paura che il regime comunista della korea del nord potesse essere messo in discussione \t ma Mao vuole mantenere uniformita ideolgoica \\
Non interviene ufficialmente ma ufficiosamente \t manda i volontari, che sono delle truppe ma senza intervento armato dichiarato \\
Questo pero capovolge le sorti \t e le truppe cinesi sconfinano nel sud korea \\
Nel 51 Truman decide di avviare delle trattative, e la guerra finisce nel 53 con la riaffermazione del 38esimo parallelo come confine

\v

Nel frattempo nel mondo ci sono altri episodi che aumentano tensione \\
Rivoluzione cubana è un fatto autonomo che succede a Cuba, ma che contribuisce alla tensione

\v

Nel 53 muore Stalin e inizia "fase del disgelo" \t sale al potere Cruscev \\
Propone un piano di riforme per migliorare condizioni dell unione sovietica, e vuole ammodernizzarla e migliorare relazioni con occidente \\
URSS deve mostrare superiorita non militarmente, ma mostrando che il sistema economico comunista è migliore rispetto a quello capitalista \\
La competizione si deve spostare sul piano economico e scientifico

\v

Negli USA diventa presidente Eisenhower \t che era generale in Normandia \\
Sembra voler consolidare la politica truman \t politica anticomunista in cui si vuole proprio lottare al comunismo, non solo contenerlo \\
La sua politica pero si mostra attenta ai segni di apertura dell unione sovietica \\
Cosi nel 55 si tiene la conferenza di Ginevra, a cui partecipano Cruscev e Eisenhower, e anche UK e Francia \t qua le grandi potenze cominciano a riparlarsi \\
Si parla quindi di clima di distensione, e poi si parlera di disgelo \t ma non è proprio cosi: tensioni importanti ci sono ancora, ma è vero che volonta di parlarsi (senza pero voler cedere) \\
Cruscev continua con la sua linea di "destalinizzazione" \t nel 56 denuncia apertamente i crimini di Stalin \t non solo cambia politica, ma condanna le sue pratiche \t si schiera contro le purge e le epurazioni \\
Questo al ?esimo congresso comunista del 56 \\
Lui scrive un rapporto segreto, che fuoriesce ed arriva ai giornali americani \t viene reso noto, e quando questo rapporto viene reso noto i partiti comunisti negli altri stati sono disorientati \t Stalin era un idolo \\
I Paesi del blocco comunista, come Polonia Ungheria etc. \t sperano che ci possa essere un comunismo meno oppressivo \\
In Polonia vengono sostituiti tutti i dirigenti di partito legati a Stalin \t si prova quindi a modificare qualcosa

\v

L'Ungheria invece vuole cambiare + radicalmente \t Nagi vuole far uscire Ungheria da patto di Varsavia e concedere liberta di stampa \\
A questo tentativo riformista, nel 56, interviene l'armata rossa che reprime nel sangue una riforma in uno stato sovrano \t nel 58 Nagi viene impiccato \\
Disgelo si perche c'è dialogo, ma entrambi i blocchi non vogliono perdere terreno \\
Questo suscita molto scalpore, anche tra i comunisti negli altri paesi \t anche in Italia si prendono le distanze un po dall'URSS \\
La repressione in Ungheria dimostra che il regime è comunque oppressivo e che non cederà nulla \t e che gli stati del blocco sovietico sono in realta stati satellite \t dipendono totalmente dalla volonta di Mosca \\
Non sono stati autonomi che condividono linea, ma subordinati \\
La primavera di Praga del 68 sara ancora + grave

\ss[Cecoslovacchia]
Dubceck vuole riformare il comunismo ceco \t "socialismo dal volto umano" era un socialismo che avesse un'economia + libera, ma soprattutto voleva un comunismo in un solo paese \t voleva una forma di comunismo che appartenesse alla sua nazione e che non fosse la copia di quello sovietico \\
URSS si preoccupa \t nuova strada del socialismo metteva in discussione URSS \\
Cosi nel ?? carri armati entrano a Praga, che segna la fine della "Primavera di praga" = periodo di tentata liberalizzazione e riforme \\
L'URSS giustifica questo attraverso la teoria della sovranita limitata \t viene inviata una lettera al comitato centrale del partito comunista cecoslovacco, in cui si dice che nessun paese del blocco sovietico è libero di cambiare sistema politico ed economico \\
Se avesse fatto questo, si sarebbe incrinata la sicurezza di tutti gli altri paesi del blocco \t la sovranita di ogni paese comunista è subordinata agli interessi dell'URSS \\
La teoria della sovranita limitata è anche chiamata "teoria Breznev", che è colui che aveva mandato la lettera \t lui nel 64 aveva sostituito URSS

\v

Nel 64 muore quindi Cruscev e arriva Breznev

\ss[Kennedy]
Nel frattempo viene eletto Kennedy, il + giovane presidente e il primo di religione cattolica (origine irlandese) \\
Ha grande popolarita per la sua storia e per il suo carisma \t ma anche non piace molto \\
Lui si impegna a rinnovare la societa americana e a portare avanti una politica di distensione internazionale \\
Nel suo programma elettorale parla di "nuova frontiera" (nel 1960 a Los Angeles, prima di essere eletto) con cui indica il nuovo scopo dell'america \\
Prima si era andati ad ovest \t ora bisogna sorpassare frontiera culturale, fare passo avanti \t parla di pace, educazione, pregiudizi non debellati \\
Non piace ad alcuni perche questo programma di riforma dava attenzione ai neri \t e la questione razziale era enorme \\
Aumenta spesa sociale, supporta ceti + bassi e finanza ricerca spaziale

\v

Sul piano internazionale Kennedy gestisce situzioni di grande tensione \t prima di tutto il problema di Berlino \\
Sovietici volevano trasformare Berlino ovest in città libera, nel senso che non era + di nessuno \t invece gli americani ritenevano che quella parte di berlino doveva fare parte della germania federale \\
Nel 61 incontra Kruschev, ma non si conclude nulla \t pero si parlano \\
I sovietici rispondo al fallimento di questo colloquio con la costruzione del muro \t nel 61 viene costruito un muro che divide berlino in due settori \\
Di voleva infatti mettere fine alle fughe delle persone che scappavano dalla zona sovietica \\
Quando il giorno dopo c'è il muro, non si fa nulla \t e rimane li, perche la spartizione del mondo in due zone di influenza (che non era ufficiale) piaceva a tutti \\
Andava bene a occidente e oriente \t andava bene anche limite fisico, cosi ognuno sta separato \t il muro diventa il simbolo della guerra fredda

\v

Poi Kennedy deve affrontare anche Cuba nel 61 \t nel 59 era andato al potere a Cuba Fidel Castro, che guida governo marxista \\
Dopo la rivoluzione del 59, vengono fatte pressioni sui presidenti usa per eliminare questo pericolo cosi vicino \t e si fa pressione anche su Kennedy \\
Kennedy fa quindi spedizione di esuli cubani che insieme ai servizi segreti dovevano sbarcare nella Baia dei Porci per fare rivoluzione \\
Sbarcano sulla baia, e si scontrano con l'esercito di Fidel Castro \t e ovviamente hanno la peggio \\
Questa sconfitta pesa sulla popolarita di Kennedy \t la reazione di Kennedy è embargo totale verso Cuba \t si vuole riscattare la sua credibilita \\
Divieto di commerciare con cuba dura 50 anni \t e avra conseguenze molto gravi per Cuba, perche ne blocca lo sviluppo economico \\
Fidel Castro teme quindi un invasione americana, e si avvicna all URSS \t che ottiene possibilita di creare base missilistica nucleare li \\
Kennedy decreta un blocco navale totale dell'isola, URSS smantella basi missilistiche, in giorni in cui si era sull'orlo della guerra \\
Dopo questo incidente nel 63 Kennedy firma con Cruscev che limita per prima volta esperimenti atomici \t ma solo nell'atmosfera, mentre limitati esperimenti nel sottosuolo
\end{document}

